\documentclass[11pt]{article}

\usepackage[margin=1.1in]{geometry}
\usepackage{amsmath,amssymb,amsthm}
\usepackage{mathtools}
\usepackage{booktabs}
\usepackage{array}
\usepackage{microtype}
\usepackage{enumitem}
\usepackage[authoryear,round]{natbib}
\usepackage[hidelinks]{hyperref}
\usepackage{cleveref}

% ---------------------------------------------------------------------------
% Theorem environments (shared by all sections)
% ---------------------------------------------------------------------------
\theoremstyle{plain}
\newtheorem{theorem}{Theorem}[section]
\newtheorem{lemma}[theorem]{Lemma}
\newtheorem{proposition}[theorem]{Proposition}
\newtheorem{corollary}[theorem]{Corollary}
\theoremstyle{definition}
\newtheorem{definition}[theorem]{Definition}
\newtheorem{premise}{Premise}
\renewcommand{\thepremise}{P\arabic{premise}}
\theoremstyle{remark}
\newtheorem{remark}[theorem]{Remark}
\newtheorem{example}[theorem]{Example}

\crefname{premise}{premise}{premises}
\Crefname{premise}{Premise}{Premises}

% ---------------------------------------------------------------------------
% Shared notation (see Section 2). Sections MUST use these macros.
% ---------------------------------------------------------------------------
\newcommand{\Sig}{\Sigma}                       % finite alphabet
\newcommand{\Str}{\Sigma^{*}}                   % finite strings
\newcommand{\Seq}{\Sigma^{\infty}}              % infinite sequences
\newcommand{\N}{\mathbb{N}}
\newcommand{\R}{\mathbb{R}}
\newcommand{\rec}{R}                            % an observation record
\newcommand{\dom}{\operatorname{dom}}
\newcommand{\ph}[1]{\varphi_{#1}}               % acceptable numbering
\newcommand{\Cons}{\operatorname{Cons}}         % models consistent with a record
\newcommand{\Sel}{\operatorname{Sel}}           % selection under a constraint
\newcommand{\K}{K}                              % prefix Kolmogorov complexity
\newcommand{\KU}[1]{K_{#1}}                     % complexity relative to machine #1
\newcommand{\Univ}{\mathbf{M}}                  % universal (Solomonoff) semimeasure
\newcommand{\conv}{\mathord{\downarrow}}        % converges
\newcommand{\divg}{\mathord{\uparrow}}          % diverges

\title{Algorithmic Agnosticism:\\
Why No Theory of Everything Can Be Known to Be Final}
\author{Christopher S. Taylor}
\date{Version 1.0 --- preprint, not peer-reviewed --- November 2026\\[0.5ex]
\small Licensed under CC BY 4.0 --- \url{https://github.com/extrafox/algorithmic-agnosticism}}

\begin{document}
\maketitle

\begin{abstract}
Physicists hope for a final theory of everything. We argue that no such
theory could be known to be final, where knowing finality means that the
observations rule out every consistent rival. We assume that every observation is
a finite record and every usable method for choosing theories is computable. From learning theory and algorithmic information
theory we show three things. First, no method of any kind can certify a theory as
final: every finite record fits infinitely many computable theories that disagree about
the unobserved. Second, no usable method is guaranteed to settle on a final theory:
each record fits some possible world in which the method never settles on the right
theory, or, if it uses chance, settles with arbitrarily small probability. Third, even
an agent with unlimited data and computation cannot separate formulations with
identical predictions or tell apart with certainty chance hypotheses that agree on
which events have probability zero. Simplicity does not decide. The data fix which
theories fit, but their lengths depend on the universal computer chosen to measure
them. Rivals that fit equally well, are nearly as short and disagree about the
unobserved always exist, and a change of computer that alters no length by more than about
their gap makes one the shortest. Weighting theories by length predicts within
as many bits of each model as it takes to describe that model, but choosing among survivors is a constraint we impose.
Compact, unified theories remain worth seeking for their testable predictions.
We call this position algorithmic agnosticism.
\end{abstract}

\noindent\textbf{Keywords:} underdetermination; scientific realism; Solomonoff induction;
minimum description length; formal learning theory; physical Church--Turing thesis;
theory of everything; agnosticism.

\tableofcontents

\section{Introduction}\label{sec:intro}

Much of fundamental physics is organized around the hope of a final theory: a single
theory of everything from which every regularity follows, and which is not itself
the approximation of something deeper \citep{weinberg1992}. Some go further and hold
that the world simply \emph{is} such a structure \citep{tegmark2008}, or that a unique
computational object underlies all possible physics \citep{wolfram-ruliad2021}. We do
not ask whether a final theory exists. We ask whether one could ever be \emph{known} to
be final (\cref{def:known}), and we argue that it could not.

The argument starts from two premises. The first is nearly definitional: every
observation, measurement or other use of a system leaves a finite record written at
finite precision (\cref{prem:record}). The second is an assumption we assert and defend
rather than prove: every system anyone can use, including every method of choosing
among theories, is computable or a computable stochastic process (\cref{prem:effective}).
It is the modest physical Church--Turing thesis \citep{piccinini2011}, extended from
devices to methods.

From these premises the main claim follows in three parts.
\begin{enumerate}[label=(C\arabic*)]
  \item \emph{No certification.} Any theory that fits every observation made so far is one
        of infinitely many that fit them equally well and disagree about what has not yet
        been observed. So no method of any kind, computable or not, can certify a theory
        as final: if a method certifies a model, deterministic or stochastic, at some
        record, a computable world that produces that record makes the certification
        wrong (\cref{thm:k1-certify}).
  \item \emph{No reliable discovery.} No usable method is guaranteed, even with any
        fixed probability, to settle on a final theory. Whatever has been observed so far,
        there is a possible world that fits all of it in which a deterministic usable
        method never settles on the right theory; for a randomized usable method and
        every $p>0$, there is such a world in which it settles with probability less than
        $p$ (\cref{thm:ud-gold,ud:rem-randomized}). A user can therefore never rule out living in such a
        world. This part rests on the second premise; for deterministic methods it needs
        only the weaker premise that every world computable at a method's own level of
        power is possible (\cref{rem:p2star,thm:ud-gold-rel}).
  \item \emph{Permanent underdetermination, even in principle.} Even an ideal agent with
        unlimited data and computing power faces two limits. (a) Formulations that make
        identical predictions cannot be separated by any method. (b) Chance hypotheses
        that agree on which events have probability zero cannot be told apart
        \emph{with certainty} by any method (\cref{thm:t5-permanent}).
\end{enumerate}
A corollary follows (C$'$): simplicity narrows the field but cannot decide among the
survivors that matter. The data fix which models fit; which of them is shortest depends
on the reference machine. Rivals that fit equally well, are within $\K(x)+O(1)$ bits of
the shortest, and disagree at an unobserved setting $x$ always exist, and a change of
machine that moves no description length by more than about that gap makes such a rival the
unique shortest (\cref{thm:con-nearmin,thm:con-gap,con:cor:promote}). The universal
mixture predicts within $\K(\nu)$ bits of every model $\nu$ (\cref{thm:pred-relative}).
Every choice among those survivors---by
simplicity, beauty, naturalness or unification---is a constraint we impose
(\cref{lem:bias}). None of this makes the search for compact,
unified theories pointless. Compact theories compress what is known and commit to
predictions about what is not, which can be tested (\cref{sec:predict}). What is
misguided is the goal of a final theory known to be final.

\paragraph{Algorithmic agnosticism.}
We call this position \emph{algorithmic agnosticism}.\footnote{The name is unrelated
to ``agnostic learning'' in computational learning theory \citep{kearns1994}.} Huxley, who coined the word, stated agnosticism as a method:
do not pretend that conclusions are certain which are not demonstrated or demonstrable
\citep[pp.~245--246]{huxley1889}. In the sense we intend, commonly glossed as neither affirming
nor denying a claim that evidence cannot settle, we neither affirm nor deny that a final
theory exists; we hold that none could be known to be one. The position differs from van
Fraassen's constructive empiricism \citep{vanfraassen1980}, which is agnostic about
unobservable entities: we are agnostic about which law generates the observations
themselves. The qualifier \emph{algorithmic} names the framework of the argument, in
which theories are programs and usable methods are algorithms. For the contrast with
Müller's \emph{algorithmic idealism} \citep{muller2020,muller2026}, see
\cref{sec:objections}.

\paragraph{What we do not claim.}
We do not claim that the world is computable, discrete, continuous, random or
deterministic, nor that it lacks a true description.

\paragraph{Precursors.}
That universal laws can in principle be refuted but never verified goes back to Hume
and was made central by Popper \citep{popper1959}. That a finite record admits arbitrarily many laws,
so that a law is worth having only if it compresses its data, goes back to Leibniz and
Weyl \citep{leibniz1686,weyl1932} and was made precise by Chaitin \citep{chaitin2005}.
Putnam proved that no computable method extrapolates every computable sequence
\citep{putnam1963}, and Gold that none identifies every computable law in the limit
\citep{gold1967}; Kelly recast scientific method in these terms
\citep{kelly1996}. Assuming an objective, computable world, Hutter formalized theories
of everything as programs selected by description length, and called such a theory
``subjective'' because it must include the observer's localization within the world
\citep[Sec.~4]{hutter2010}. He also notes that ``for every model we can construct an
alternative model that behaves identically in the past'' and predicts differently
\citep[Sec.~4]{hutter2010}, which is our \cref{thm:ud-infinite}, and he formalizes and examines a
universal theory of everything that generates all computable universes
\citep[Sec.~6]{hutter2010}; choosing among such theories by description length is
machine-relative (\cref{thm:con-machine}). Schmidhuber treats a theory of everything as a
prior over computable universes and notes that assuming it formally describable
``imposes a strong inductive bias'' \citep[abstract]{schmidhuber2000}; \cref{lem:bias} makes the
point general. Sterkenburg argues that Solomonoff prediction cannot justify a
preference for simplicity: its predictors are mixtures over effective hypotheses, so
the operative assumption is effectiveness, and the resulting ``simplicity'' depends on
an arbitrary choice of universal machine \citep{sterkenburg2016}. A strong case for
simplicity is Vitányi's: ideal MDL is optimal when the data are typical for the true
hypothesis, and algorithmic statistics separates a record's meaningful information from
its noise \citep{vitanyi-li2000,gacs2001}; we answer it in \cref{sec:simplicity,sec:objections}.
Hawking argued from
G\"odel's theorem that a final theory may be out of reach \citep{hawking2002}, and
Faizal and colleagues argue from incompleteness that a theory of everything would require
non-algorithmic understanding \citep{faizal2025}. Kiefer concludes from G\"odel's theorems that,
unless space-time is fundamentally discrete, ``we can never decide whether a given
theory is the final one or not'' \citep[abstract]{kiefer2024}. Our route to a similar conclusion
runs through underdetermination by finite records and the limits of learning, not
through undecidability within the theory, so it does not depend on the theory's
arithmetic strength or on whether space-time is discrete. Hossenfelder argues that the
appeal to beauty and naturalness has misled the search for a final theory \citep{hossenfelder2018}.

\paragraph{Contributions.}
Little of the mathematics is new; the contribution is the assembly, aimed at a specific
target. Specifically:
(i) we show that the core of the claim, no certification (C1) and permanent
underdetermination (C3), holds for every method whatsoever;
(ii) we state Gold's theorem in a form in which, whatever record has been observed, some
world that fits it is one the method never gets right, which answers the reply that a
method need only succeed in \emph{our} world;
(iii) we make the machine-relativity of simplicity quantitative for theories fitted to
finite data (\cref{thm:con-machine}), rather than leaving it as a remark about arbitrary
constants;
(iv) we isolate the underdetermination that is permanent even for an ideal agent
(\cref{thm:t5-permanent}); and
(v) we reach a conclusion like Kiefer's (above) from finite records and the limits of
learning alone.

\paragraph{Outline.}
\Cref{sec:framework} fixes definitions and states and defends the premises.
\Cref{sec:certify} shows that no record certifies a final theory, and
\cref{sec:discover} that no usable method reliably finds one.
\Cref{sec:simplicity} shows that simplicity does not single one out, and
\cref{sec:ideal} what even an ideal agent cannot know.
\Cref{sec:predict} explains what the search is nonetheless good for.
\Cref{sec:objections} answers objections, and \cref{sec:conclusion} concludes.

\section{Framework and premises}\label{sec:framework}

\subsection{Records, models, worlds}

Fix a finite alphabet $\Sig$ with $|\Sig|\ge 2$. Write $\Str$ for the finite strings
over $\Sig$, $\Seq$ for the infinite sequences, and $u\sqsubseteq v$ for the prefix
order on strings and sequences.

\begin{definition}[Record]\label{def:record}
An \emph{observation} is a pair $(x,y)\in\Str\times\Str$: a setting $x$ and a
recorded outcome $y$. A \emph{record} is a finite sequence
$\big((x_0,y_0),\dots,(x_{n-1},y_{n-1})\big)$ of observations; repetitions are
allowed, so repeated trials keep their order and multiplicity. Records grow by
extension. We write $\rec$ also for the underlying set of observations when order
does not matter, and $\dom(\rec)=\{x_i\}_{i<n}$. A record is \emph{functional} if no
setting occurs with two distinct outcomes. Fixing a prefix-free pairing code, a
record is encoded as a single string $r\in\Str$, and extension of records is prefix
extension of their codes.
\end{definition}

Fix an acceptable (G\"odel) numbering $(\ph{e})_{e\in\N}$ of the partial computable
functions $\Str\rightharpoonup\Str$ \citep{rogers1967}; write $\ph{e}(x)\conv$ when
the computation halts.

\begin{definition}[Model]\label{def:model}\label{def:probmodel}
(a) A \emph{deterministic model} is an index $e\in\N$. It is \emph{consistent} with a
functional record $\rec$ if $\ph{e}(x)\conv=y$ for every $(x,y)\in\rec$; we write
$\Cons(\rec)$ for the set of such $e$. Two models \emph{diverge off} $\rec$ if they
give different outputs at some setting not in $\dom(\rec)$.
(b) For a fixed sequence of settings, fix an effective enumeration
$\nu_1,\nu_2,\dots$ of the lower semicomputable semimeasures on outcome sequences
\citep{li-vitanyi2019}. A \emph{stochastic model} is an index $i\ge1$; $\nu_i(r)$ is the
probability it assigns to outcomes beginning with $r$, and it is consistent with a
record with outcome string $r$ if $\nu_i(r)>0$.
(c) Two models are \emph{extensionally equal} if they have the same extension:
$\ph{e}=\ph{e'}$ for deterministic models, $\nu_i=\nu_j$ for stochastic ones.
\end{definition}

\begin{definition}[World]\label{def:world}
A \emph{world} is whatever produces the records. A world is \emph{computable} if its
record-producing behavior is given by a total computable function $f:\Str\to\Str$
(deterministic) or a computable measure $\mu$ (stochastic). We make no claim that the
actual world is computable; computable worlds are the worlds that some usable model
could describe exactly.
\end{definition}

\begin{definition}[Final model; theory of everything]\label{def:final}
A model is \emph{final} for a world if it agrees with every record the world can
produce: $\ph{e}=f$ in the deterministic case, $\nu_i=\mu$ in the stochastic case. A
\emph{theory of everything} is any model offered as final.
\end{definition}

\begin{definition}[Method; certification; identification]\label{def:learner}
A \emph{method} is a function from (codes of) records to models; a \emph{randomized}
method also reads bits from a random source. A method \emph{certifies} a model at a
record $\rec$ if on $\rec$ it outputs that model together with the verdict ``final''.
Let $x_0,x_1,\dots$ be the length-lexicographic enumeration of the settings. The
\emph{canonical presentation} of a total $f$ is the sequence of records
$\rec_{f\restriction n}=\big((x_0,f(x_0)),\dots,(x_{n-1},f(x_{n-1}))\big)$, $n\in\N$.
A method that has already observed a record $\rec$ consistent with $f$ is presented
with $f$ \emph{after} $\rec$: it receives the records $\rec\frown\rec_{f\restriction n}$,
$n\in\N$ (for empty $\rec$, the canonical presentation). It \emph{identifies $f$ in
the limit} after $\rec$ if its outputs on this presentation are eventually constant at
an index of $f$ \citep{gold1967}; it does not have to announce when this has happened.
A method that is undefined on some record of the presentation does not identify $f$.
A randomized method identifies $f$ \emph{with probability $p$} if the set of random
bit sequences on which it identifies $f$ has probability $p$. A method is
\emph{guaranteed} to find a final theory after $\rec$ if it identifies in the limit,
after $\rec$, every total computable world consistent with $\rec$; a randomized method
is guaranteed with probability $p$ if it identifies each such world with
probability at least $p$.
\end{definition}

\begin{definition}[Known to be final]\label{def:known}
A theory $T$ is \emph{known to be final}, in the sense used here, on a record $\rec$ if
$\rec$ rules out every rival consistent with it: $T$ is final (\cref{def:final}) for every computable
world that can produce $\rec$: for a deterministic theory, every total computable $f$ consistent with
$\rec$; for a stochastic theory, every computable measure $\mu$ with $\mu(r)>0$, where
$r$ is the outcome string of $\rec$.
Equivalently, a method that certifies $T$ at $\rec$ is right in every such world.
\end{definition}

We take \cref{def:known} as the standard for the hope of a final theory: finality recognized
from the evidence. It counts only computable
worlds as rivals, so it is weaker than a standard that counts every world, and a result
that defeats it defeats every stronger one.

\begin{definition}[Constraint; selection]\label{def:constraint}
A \emph{constraint} (or \emph{bias}) is given by a cost $c_\rec:\N\to\N\cup\{\infty\}$
on models, possibly depending on the record. Its \emph{selection} at a functional
record $\rec$ is
\[
  \Sel_c(\rec)=\operatorname*{arg\,min}_{e\in\Cons(\rec)} c_\rec(e).
\]
A constraint is \emph{imposed} if the record does not determine it: other constraints,
equally consistent with the record, select differently on the same record, and nothing
in the record decides between them.
\end{definition}

``Which model is final'' can be asked at three levels: of \emph{programs} (indices),
of \emph{extensions} (the behavior $\ph{e}$), and of \emph{chance hypotheses}
(measures).

\subsection{Premises}

\begin{premise}[Finite record]\label{prem:record}
Every use of any system by any human or machine yields a record in the sense of
\cref{def:record}.
\end{premise}

\Cref{prem:record} is close to definitional. Uses take finite time; instruments have
finite resolution; outcomes are written down, stored or acted upon in finite memory.
Even an instrument reading a continuous quantity reports a finite-precision value.
\Cref{prem:record} is a claim about \emph{interfaces}, not about the systems behind
them: it says nothing about whether the world is discrete, continuous, computable or
random.
\emph{Refuted by:} a use of some system whose outcome is stored or acted on as an
infinite object, for instance an instrument whose reading is available as an exact
real number.

\begin{premise}[Usable systems are effective]\label{prem:effective}
Every system that any human or machine can use to produce outputs from inputs---in
particular every method by which records are turned into predictions or model
choices---is effective: its input--output behavior is computable or, if it is
stochastic, is described by a computable probability measure on outputs.
\end{premise}

\Cref{prem:effective} is an assumption; it does not follow from \cref{prem:record}.
It is Piccinini's \emph{modest} physical Church--Turing thesis
\citep{piccinini2011,piccinini2015}---every function computable by a process an
observer can \emph{use} is Turing-computable---whose usability conditions we adopt:
readable inputs and outputs, a rule specifiable independently of the process,
repeatability, settability, physical constructibility and reliability
\citep{piccinini2015}, a list he does not intend to be exhaustive. We extend it
in scope only: through the measure clause to stochastic systems, and to every usable
method, including methods of choosing models.
A stochastic method in the sense of \cref{prem:effective} is a computable machine with
access to fair coin flips, the model of \citet{pitt1989}; every computable measure on
outputs is realized by such a machine.

\emph{Refuted by:} a usable hypercomputer, that is, a system whose non-computable
behavior can be specified, relied on and put to work.

\paragraph{Defense.} The three standard candidates for a usable hypercomputer fail.

\emph{(a) Continuous theories.} A theory stated over the reals is used by computing
its consequences to finite precision, and whatever a theory yields in that way is
computable in the sense of computable analysis \citep{weihrauch2000}. Some
continuous theories have non-computable values---the wave equation has computable
initial data whose unique (weak) solution at a later time is not computable
\citep{pourel-richards1981}---and recurrent networks with real-valued weights decide
non-recursive languages \citep{siegelmann-sontag1994}. But such values cannot be
used: no finite record can read them out. In the network case the non-computability
comes out only where it has been put in, through non-computable weights
\citep{davis2004}.

\emph{(b) Relativistic hypercomputers.} In Malament--Hogarth spacetimes, for example
that of a rotating (Kerr) black hole, an observer could in principle receive, in finite proper time, the result of an infinite
computation \citep{hogarth1992,etesi-nemeti2002}. But the observer cannot rule out
that a ``does not halt'' verdict reflects spurious signals or a failure of the device
rather than an infinite computation, so the output cannot be relied on; Earman and
Norton, who make this point, regard such computation as an open question rather than
an impossibility \citep{earman-norton1993}.

\emph{(c) Quantum randomness.} A random source can produce a non-computable sequence,
but it cannot be steered to compute anything specified in advance.

\begin{theorem}[Randomness adds no power to compute functions;
  {\citealp[Thms.~2--3]{deleeuw1956}}]\label{thm:q-powerless}
If $p$ is a computable real and a probabilistic machine with a $p$-biased coin
enumerates a set with positive probability, the set is computably enumerable; if it
outputs $f(n)$ with probability greater than $\tfrac12$ on every input $n$, then $f$
is computable.
\end{theorem}

\begin{proposition}[Born statistics are computable]\label{prop:q-born-computable}
Let a finite-dimensional quantum system have a computable initial state $\psi$, a
computable unitary $U$ and a measurement with computable projectors $\{P_k\}$. Then
the outcome probabilities $\lVert P_kU\psi\rVert^2$ \citep{born1926} are computable
reals, and the outcomes of repeated independent measurements form a computable
measure.
\end{proposition}

\begin{proof}[Proof sketch]
Matrix products, exponentials of computable matrices and norms are computable
operations on computable reals \citep{weihrauch2000,pourel-richards1989}; a product
of computable probabilities defines a computable product measure.
\end{proof}

A quantum random number generator is therefore a computable stochastic process in
the sense of \cref{prem:effective}: what can be relied on about it is a computable
measure, and by \cref{thm:q-powerless} its randomness computes nothing new. A coin
with a non-computable bias would act as an oracle, but no finite record could show
that its bias is non-computable, so it could not be relied on as one.

\subsection{The argument in outline}\label{sec:outline}

\Cref{tab:outline} is the map of the argument. It lists every step in the order of the
sections that prove it, with the level at which each holds.

\begin{table}[!htbp]
\centering
\small
\begin{tabular}{@{}l>{\raggedright\arraybackslash}p{0.55\linewidth}%
>{\raggedright\arraybackslash}p{0.13\linewidth}%
>{\raggedright\arraybackslash}p{0.17\linewidth}@{}}
\toprule
Step & Statement & Status & Level \\
\midrule
Known & To know that a theory is final is for the record to rule out every
  computable rival (\cref{def:known}) & definition & --- \\
\ref{prem:record} & Every use yields a finite record & premise & --- \\
\ref{prem:effective} & Every usable system, every method included, is computable or
  a computable stochastic process & premise & --- \\
T1 & Every record has infinitely many consistent models, deterministic or stochastic,
  that diverge on unobserved settings (\cref{thm:ud-infinite}) & theorem &
  programs, extensions, chances \\
K1 & If a method certifies a model at a record, some computable world producing that
  record makes the certification wrong (\cref{thm:k1-certify}) & corollary &
  extensions, chances \\
T2 & For every deterministic computable method and every record, some computable world
  consistent with the record defeats it; a randomized one succeeds with no fixed
  positive probability (\cref{thm:ud-gold}; for active methods, \cref{rem:active}; at any oracle level,
  \cref{thm:ud-gold-rel})
  & theorem & extensions \\
P2$^{*}$ & Every usable method faces all worlds computable at its own level
  (\cref{rem:p2star}) & premise, weaker than P2 & --- \\
L & Every consistent method is a bias that the record does not determine (\cref{lem:bias}) &
  lemma & --- \\
T3 & Which consistent model is shortest depends on the reference machine: a near-tie of
  $\delta$ bits is reversed by a change of machine that moves no description length by
  more than $\delta$, and rivals within about $\K(x)$ bits of the shortest disagree at
  any unobserved setting $x$
  (\cref{thm:con-machine,thm:con-gap,thm:con-nearmin,con:cor:promote}) & theorem & extensions \\
T5 & No method separates formulations with identical predictions, or mutually
  absolutely continuous chance hypotheses with probability 1 under both (\cref{thm:t5-permanent}) & theorem &
  programs, chances \\
T4 & The universal predictor is within $K(\nu)$ bits of every model $\nu$; no
  computable predictor is (\cref{thm:pred-relative,pred:prop:no-computable}) &
  theorem & extensions, chances \\
\midrule
C1 & No method of any kind can certify a theory, deterministic or stochastic, as final
  & from T1, K1 & \\
C2 & No usable method is guaranteed, even with any fixed probability, to settle on a
  final theory & from P2 (P2$^{*}$ for deterministic methods), T2, L & \\
C3 & Formulations with identical predictions, and chance hypotheses that agree on
  which events have probability zero, are never told apart with certainty & from T5
  (T1(b) illustrates) & \\
C$'$ & Simplicity narrows the field but cannot decide among near-tied survivors that
  disagree about the unobserved; choosing among them is an imposed constraint;
  compression still predicts & from L, T3, T4; P2 to read T4's benchmark as unusable & \\
\bottomrule
\end{tabular}
\caption{The steps of the argument in section order. ``Level'' says whether a step
concerns programs (indices), extensions (behavior) or chance hypotheses (measures).
With \cref{def:known}, C1 is the main claim that no theory can be known to be final.}
\label{tab:outline}
\end{table}

\noindent C2 uses \cref{prem:effective}, and for deterministic methods only its weaker form
P2$^{*}$. C$'$ uses \cref{prem:effective} only to read the benchmark of T4 as unusable.

\section{No record certifies a final theory}\label{sec:certify}

This section proves that no finite record singles out a final theory, and that no
method of any kind can therefore certify one. The results concern records alone, and
they hold whatever produced the records.

\begin{theorem}[Infinite underdetermination]\label{thm:ud-infinite}
Let $\rec$ be any finite functional record.
\begin{enumerate}[label=(\alph*)]
  \item\label{ud:item-a} For every $x\notin\dom(\rec)$ and every $y\in\Str$ there is a
        model $e\in\Cons(\rec)$ with $\ph{e}$ total and $\ph{e}(x)=y$. Consequently
        $\Cons(\rec)$ contains an infinite family of total models any two of which
        differ at \emph{every} $x\notin\dom(\rec)$.
  \item\label{ud:item-b} For every $e\in\Cons(\rec)$ there are infinitely many
        $e'\ne e$ with $\ph{e'}=\ph{e}$; all of them lie in $\Cons(\rec)$.
  \item\label{ud:item-c} Let $r\in\Str$ be the outcome string of any finite record,
        functional or not (\cref{def:model}). There are infinitely many computable
        probability measures $\nu$ on $\Seq$ with $\nu(r)>0$ whose predictions after
        $r$ pairwise differ: $\nu(0\mid r)\ne\nu'(0\mid r)$ for $\nu\ne\nu'$, where
        $\nu(a\mid r)=\nu(ra)/\nu(r)$.
\end{enumerate}
\end{theorem}

\begin{proof}
\ref{ud:item-a} Since $\dom(\rec)$ is finite, the function
\[
  f_{x,y}(z)=
  \begin{cases}
    y' & \text{if } (z,y')\in\rec,\\
    y  & \text{if } z=x,\\
    \varepsilon & \text{otherwise}
  \end{cases}
\]
differs from the constant function $z\mapsto\varepsilon$ on finitely many arguments,
so it is total computable; any index $e$ for it lies in $\Cons(\rec)$ and has
$\ph{e}(x)=y$. For the stronger family, let $g_y(z)$ be $y'$ if $(z,y')\in\rec$ and
$y$ otherwise. Each $g_y$ is total computable and consistent with $\rec$, and for
$y\ne y'$ the models $g_y,g_{y'}$ differ at every $z\notin\dom(\rec)$.

\ref{ud:item-b} This is the padding lemma for acceptable numberings: there is a
computable injection $p:\N\times\N\to\N$ with $\ph{p(e,n)}=\ph{e}$ for all $e,n$
\citep[ch.~1]{rogers1967}. Membership in $\Cons(\rec)$ depends only on $\ph{e}$.

\ref{ud:item-c} Let $k=|\Sig|$ and $m=|r|$. For each rational $q\in(0,1)$ let $\nu_q$ make
the first $m$ symbols independent and uniform on $\Sig$, and every later symbol
independently $0$ with probability $q$ and each other symbol with probability
$(1-q)/(k-1)$. Each $\nu_q$ is computable, $\nu_q(r)=k^{-m}>0$, and
$\nu_q(0\mid r)=q$.
\end{proof}

The construction is folklore: Kukla's recipe for a rival to any
theory $T$---the theory $T''$ that agrees with $T$ wherever the world is observed and
with some other theory elsewhere \citep{kukla1993}---is the family $g_y$ made precise.
Part~\ref{ud:item-b} gives models that agree everywhere and differ only as programs;
we return to it in \cref{sec:ideal}.

\begin{corollary}[No certification]\label{thm:k1-certify}
Let $M$ be any method whatever: computable or not, deterministic or randomized.
\begin{enumerate}[label=(\alph*)]
  \item\label{k1:item-a} If at a finite functional record $\rec$ the method $M$
        certifies a deterministic model $e$ as final (\cref{def:final}) with positive
        probability, then there is a total computable world $f$ consistent with $\rec$,
        with $\ph{e}\ne f$, in which $M$ receives $\rec$ and certifies $e$ with the same
        probability.
  \item\label{k1:item-b} If at a record with outcome string $r$ the method $M$
        certifies a stochastic model $i$ as final with positive probability, then
        there is a computable measure $\mu\ne\nu_i$ with $\mu(r)>0$; in the world $\mu$,
        $M$ certifies $i$ wrongly with positive probability.
\end{enumerate}
Hence every method that certifies at some record is wrong on some computable world
that produces that record. A certification can be correct only relative to a
restricted class of worlds, one that excludes the rivals of \cref{thm:ud-infinite}.
\end{corollary}

\begin{proof}
\ref{k1:item-a} Pick $x\notin\dom(\rec)$. If $\ph{e}(x)$ is undefined, every total
$g_y$ of \cref{thm:ud-infinite}\ref{ud:item-a} is a world consistent with $\rec$ in
which $e$ is not final. Otherwise choose $y\ne\ph{e}(x)$; the world $g_y$ is
consistent with $\rec$ and differs from $\ph{e}$ at $x$. In either case the world
produces the record $\rec$, and $M$'s output depends only on its input (and, if $M$
is randomized, on its own random bits), so $M$ certifies $e$ there with the same
probability, wrongly.

\ref{k1:item-b} Take $\mu=\nu_q$ from the proof of
\cref{thm:ud-infinite}\ref{ud:item-c}, with $q\ne\nu_i(0\mid r)$ if $\nu_i(r)>0$ and
$q$ arbitrary if $\nu_i(r)=0$. Then $\mu(r)>0$ and $\mu\ne\nu_i$. In the world $\mu$ the
outcome string $r$ occurs with probability $\mu(r)>0$, and on receiving it $M$
certifies $i$ with the same positive probability as before, wrongly.
\end{proof}

The restriction to a class of worlds is a bias in the sense of \cref{def:constraint}: it is
not supplied by the record.

\begin{remark}[Each rival is refuted eventually]\label{ud:rem-transient}
Every model that gives a different output from a world $f$ at some setting is
refuted after finitely many observations of the canonical presentation of $f$. (A
model that is undefined at some setting leaves $\Cons$ once that setting is
observed, though no finite computation exhibits its failure.) So each rival constructed in
\cref{thm:ud-infinite}\ref{ud:item-a} is only transiently consistent with the
growing record, and one might hope that a method could at least \emph{converge} on
a final theory without ever certifying it. \Cref{sec:discover} shows that no usable
method can guarantee even this. The rivals that are never refuted---models that
agree on every possible record---are the subject of \cref{sec:ideal}.
\end{remark}

\section{No usable method finds one}\label{sec:discover}

\Cref{thm:k1-certify} rules out certifying a final theory. A weaker goal remains: a
method might converge on a final theory without ever announcing that it has. This
section shows that no usable method is guaranteed, even with any fixed probability, to
settle on a final theory, and that whatever success a method has is bought by a bias
the record does not supply.

Fix the length-lexicographic enumeration $x_0,x_1,x_2,\dots$ of $\Str$, a computable
bijection $\N\to\Str$, and identify a total $f:\Str\to\Str$ with the sequence
$f(x_0),f(x_1),\dots$. For a finite sequence $\sigma=(y_0,\dots,y_{n-1})$ write
$\rec_\sigma=\big((x_0,y_0),\dots,(x_{n-1},y_{n-1})\big)$, write $\sigma\tau$ for
concatenation, $\sigma\sqsubseteq\tau$ for the prefix order, and
$f\restriction n=(f(x_0),\dots,f(x_{n-1}))$. For a record $\rec$ already in hand,
$\rec\frown\rec_\tau$ is the record $\rec$ followed by the observations of $\rec_\tau$;
the presentation of $f$ after $\rec$ is the chain
$\rec\frown\rec_{f\restriction 0}\sqsubseteq\rec\frown\rec_{f\restriction 1}
\sqsubseteq\cdots$ (\cref{def:learner}). Fix two distinct symbols $0,1\in\Sig$.

Any finite functional record $\rec$ determines a finite sequence $\sigma_\rec$: take
the least $n$ with $\dom(\rec)\subseteq\{x_0,\dots,x_{n-1}\}$ and set the $i$th entry to
the outcome $\rec$ records at $x_i$, or to $0$ if it records none. Every total $f$
extending $\sigma_\rec$ is consistent with $\rec$.

\begin{theorem}[T2: Non-identifiability from any record; after \citet{putnam1963,gold1967}]\label{thm:ud-gold}
Let $L$ be a computable method (\cref{def:learner}), possibly partial, and $\rec$ a
finite functional record. There is a total computable world $f$,
consistent with $\rec$, that $L$ does not identify in the limit after $\rec$.
\end{theorem}

\begin{proof}
We build finite sequences $\sigma_0\sqsubseteq\sigma_1\sqsubseteq\cdots$ with
$\sigma_0=\sigma_\rec$. Given $\sigma_s$, search, dovetailing over $k$ and computation
time, for a $k$ such that $h=L(\rec\frown\rec_{\sigma_s0^k})$ is defined and
\[
  \ph{h}\bigl(x_{|\sigma_s|+k}\bigr)\conv = 0,
\]
that is, the current conjecture predicts that the record continues with $0$. If the
search succeeds, set $\sigma_{s+1}=\sigma_s0^k1$.

Suppose the search at some stage $s$ never succeeds. Then $L$ does not identify
$g=\sigma_s0^\infty$ after $\rec$. If it did, it would be defined on every record of
the presentation of $g$ after $\rec$, and from some point on its conjecture would be a
fixed index for $g$; that index outputs $0$ at every argument beyond $\sigma_s$, so the
search would succeed. The world $g$ is total computable and extends $\sigma_\rec$, so it
is consistent with $\rec$.

Otherwise every stage succeeds, and $f=\bigcup_s\sigma_s$ is total computable and
extends $\sigma_\rec$, so it is consistent with $\rec$. The conjecture $h_s$ found at
stage $s$ is the output of $L$ on a record of the presentation of $f$ after $\rec$,
since $\sigma_s0^k\sqsubseteq f$, and it predicts $0$ where $f$ has $1$, so
$\ph{h_s}\ne f$. These records grow with $s$. If $L$ converged on $f$ to some $h^*$, then
$h_s=h^*$ for all large $s$, so $\ph{h^*}\ne f$.
\end{proof}

A method that does well on the simplest worlds (here, those that are eventually
constant) for that very reason never settles on the right theory in another computable
world, and that world can be chosen to agree with any record already in hand. The world
is not exotic. The construction is uniform in $L$ and $\rec$. If every stage succeeds,
$\K(f)\le\K(L,\rec)+O(1)\le\K(L)+\K(\rec)+O(1)$, where $\K(L)$ is the complexity of an
index for $L$. If stage $s$ stalls, the world is eventually constant and
$\K(g)\le\K(L,\rec)+\K(s)+O(1)$. Either way its complexity is set by the method, the
record and at most one stage number, and its records are compressible, not atypical or
non-stochastic in the sense of algorithmic statistics \citep{gacs2001}.

\begin{remark}[Randomized methods]\label{ud:rem-randomized}
A method may also flip a coin of computable bias. The theorem covers such methods in
the following form. A probabilistic method
that identifies every computable world with probability greater than $1/2$ can be
replaced by a deterministic computable method that identifies the same worlds
\citep[Thm.~3.21 with $n=1$]{pitt1989}, so by \cref{thm:ud-gold} no probabilistic method identifies every
computable world with probability greater than $1/2$. Success with a smaller probability, bounded below by
a fixed $p>0$ for every world, corresponds to identification by a finite team of
deterministic methods, and the class of all computable worlds is not identifiable by
any finite team either \citetext{\citealp[Thm.~3.36]{pitt1989}; \citealp[Thm.~3.1]{smith1982}; \citealp{pitt-smith1988}}. Both facts hold
relative to any record: the class of computable worlds extending $\sigma_\rec$ is not
finite-team identifiable, since a team for it would yield one for all computable worlds
via $g\mapsto\sigma_\rec g$ and the corresponding computable shift of indices. So, given
any record, no usable randomized method that receives the presentation identifies every
computable world consistent with it with any fixed positive probability. This is a
uniform statement: it does not say that the method fails almost surely in some single
world, and a randomized method can identify every computable world with \emph{some}
positive probability, for instance by guessing an index at random and keeping it.
\end{remark}

\begin{remark}[Active methods]\label{rem:active}
A method may also choose which settings to observe. An \emph{active} method reads the
record so far, and possibly random bits, and outputs a conjecture and a next setting to
query; the world answers the query, and the answer extends the record. Active methods
gain nothing here. Given a computable active method $M$ and a record $\rec$, let $L_M$
be the method that receives the presentation after $\rec$ and, on the record
$\rec\frown\rec_{f\restriction n}$, simulates $M$ from $\rec$ for $n$ computation steps
with the same random bits, answering each query from the record and halting the
simulation at the first query the record does not yet answer; $L_M$ outputs the last
conjecture of $M$ reached. $L_M$ is computable, and randomized exactly when $M$ is.
Every setting is eventually answered, so every step of a run of $M$ in which each
step halts is eventually reached; if on given random bits $M$ identifies $f$, its
conjectures are eventually constant at an index of $f$, and so are those of $L_M$.
Hence $L_M$ identifies every world that $M$ identifies, with at least the same
probability. Applied to $L_M$, \cref{thm:ud-gold} and \cref{ud:rem-randomized} give,
for every record, a computable world consistent with it that $M$ does not identify, or,
for randomized $M$ and every $p>0$, one that $M$ identifies with probability less than
$p$.
\end{remark}

\begin{corollary}[No record rules out failure]\label{ud:cor-every-method-fails}
Let $L$ be a usable method, possibly partial, and $\rec$ a finite functional record that
$L$ has observed; $L$ may receive the presentation or choose its own queries. If $L$
is deterministic, some computable world consistent with $\rec$ \emph{defeats} it: in
that world, $L$ never settles on the right theory. If $L$ is randomized, then for every
$p>0$ there is a computable world consistent with $\rec$ that $L$ identifies with
probability less than $p$. In either case no
record tells $L$'s user that the world is not such a world.
\end{corollary}

\begin{proof}
By \cref{prem:effective} $L$ is computable or a computable probabilistic method;
apply \cref{thm:ud-gold} or \cref{ud:rem-randomized}, and for active methods
\cref{rem:active}, to $\rec$.
\end{proof}

\paragraph{The limit does not depend on the level of computing power.}
\Cref{thm:ud-gold} is not an artifact of ordinary computers. Fix any set $A\subseteq\N$
as an oracle, and let $(\ph{e}^A)_{e\in\N}$ be an acceptable numbering of the functions
partial computable relative to $A$, which exists by the relativized universal-function
and $s$-$m$-$n$ theorems \citep[ch.~9]{rogers1967}; the proof below uses only that the universal
function is $A$-computable. Call a model at level $A$ an index $e$ of
$\ph{e}^A$, and a world at level $A$ a total $A$-computable function; identification
at level $A$ is defined as in \cref{def:learner}, with $\ph{e}^A$ in place of
$\ph{e}$.

\begin{theorem}[T2 at every level]\label{thm:ud-gold-rel}
Let $A\subseteq\N$, let $L$ be an $A$-computable method, possibly partial, that outputs
models at level $A$, that is, indices in $(\ph{e}^A)$ rather than ordinary indices, and
let $\rec$ be a finite functional record. There is a total $A$-computable world $f$,
consistent with $\rec$, that $L$ does not identify at level $A$ after $\rec$.
\end{theorem}

\begin{proof}
The proof of \cref{thm:ud-gold} goes through with $\ph{}^A$ in place of $\ph{}$. Given
$\sigma_s$ (with $\sigma_0=\sigma_\rec$), search for $k$ such that
$h=L(\rec\frown\rec_{\sigma_s0^k})$ is defined and $\ph{h}^A(x_{|\sigma_s|+k})\conv=0$.
The computations of $L$ and of $\ph{h}^A$ can be simulated with access to $A$, so the
search, dovetailed over $k$ and computation time, is $A$-effective. If it never
succeeds at stage $s$, then $L$ does not identify $\sigma_s0^\infty$ at level $A$, for
the same reason as before, and that world is computable, hence $A$-computable.
Otherwise set $\sigma_{s+1}=\sigma_s0^k1$. Then $f=\bigcup_s\sigma_s$ is total, extends
$\sigma_\rec$, and is $A$-computable, since every stage is effective relative to $A$.
Each conjecture $h_s$ found at stage $s$ predicts $0$ where $f$ has $1$, so if $L$
converged on $f$ to some $h^*$ then $\ph{h^*}^A\ne f$.
\end{proof}

The relativized theorem needs the worlds to be at the method's level: a method strictly
more powerful than the worlds it faces can succeed, as the ideal-agent concession below
shows. For deterministic methods, C2 therefore needs only the following premise.

\begingroup
\renewcommand{\thepremise}{P2$^{*}$}%
\renewcommand{\theHpremise}{P2star}%
\addtocounter{premise}{-1}%
\begin{premise}[Methods and worlds share a level]\label{rem:p2star}
For every usable method $L$ there is a set $A\subseteq\N$ such that $L$ is $A$-computable, its outputs are read as models at level $A$, and its guarantee must cover every total $A$-computable world consistent with the record: every such world is a possible world. A model at level $A$ is final for a world $f$ if it computes $f$ relative to $A$.
\end{premise}
\endgroup

\noindent\emph{Refuted by:} a usable method more powerful than every world its guarantee
must cover, for instance a usable oracle whose power the theories it chooses among
cannot share.

\Cref{prem:effective}, together with the guarantee of \cref{def:learner}, which covers
every total computable world consistent with the record, is the case $A=\emptyset$ of
P2$^{*}$. Under P2$^{*}$, \cref{thm:ud-gold-rel} applies at the method's own level, so
a usable hypercomputer of oracle strength, whose power the theories share, does not
escape C2. Hypercomputers not equivalent to a single oracle, such as infinite-time
Turing machines, would need their own argument. The
reductions of \citet{pitt1989} behind \cref{ud:rem-randomized} are uniform simulations
and should relativize, but we have not verified this, so P2$^{*}$ serves for
deterministic methods only.

\begin{lemma}[L: Every consistent method is a bias]\label{lem:bias}
Let $L$ be a method that outputs, on each functional record $\rec$, a model in
$\Cons(\rec)$.
\begin{enumerate}[label=(\roman*)]
  \item $L$ is a record-indexed constraint (\cref{def:constraint}): with
        $c_\rec(e)=0$ if $e=L(\rec)$ and $c_\rec(e)=1$ otherwise,
        $\Sel_c(\rec)=\{L(\rec)\}$. For every functional record $\rec$ there is a
        method $L'$ that outputs consistent models, is computable if $L$ is, and
        selects differently from $L$ at $\rec$.
  \item For a finite functional record $\rec$, let $K_L$ be the class of total
        computable worlds consistent with $\rec$ that $L$ identifies after $\rec$. If
        $L$ is computable (deterministic), some world consistent with $\rec$ lies
        outside $K_L$; and if $L$ identifies after $\rec$ the eventually-zero worlds
        consistent with $\rec$, some world lies inside $K_L$.
\end{enumerate}
\end{lemma}

\begin{proof}
(i) The first claim is immediate from the definition, since $L(\rec)\in\Cons(\rec)$ has
cost $0$ and every other model cost $1$. For the second, $\Cons(\rec)$ is infinite by
\cref{thm:ud-infinite}, so it contains some $e'\ne L(\rec)$. Let $L'$ agree with $L$ on
every record other than $\rec$, and let $L'(\rec)=e'$. Then $L'$ outputs consistent
models, it differs from $L$ only on one input, so it is computable if $L$ is, and
$\Sel_{c'}(\rec)=\{e'\}\neq\{L(\rec)\}$, where $c'$ is the constraint of $L'$ given by
(i). So the record does not determine the
constraint, which is what it means for the constraint to be imposed
(\cref{def:constraint}). For (ii), the first claim is \cref{thm:ud-gold}; for the
second, the world $\sigma_\rec0^\infty$ is eventually zero and consistent with $\rec$.
\end{proof}

So choosing a method is choosing a class $K_L$ of worlds to be right about, and no
record decides whether the actual world belongs to that class. \Cref{sec:simplicity}
examines simplicity as a way of making this choice.

\paragraph{The ideal-agent concession.}
\Cref{thm:ud-gold} binds computable methods, and the restriction is not idle. A
method with an oracle for the totality problem can output the least total index
consistent with the current record; every smaller total index computes a different
function and is eventually refuted, so this \emph{identification by enumeration}
converges to a correct model on every total computable world
\citep{putnam1963,gold1967}. An ideal agent of this
kind therefore does converge on the true behavior of any deterministic computable
world, though it never knows when it has. The limit of \cref{thm:ud-gold} is
computational.

\section{Simplicity does not single one out}\label{sec:simplicity}

The most common preference in the search for a theory of everything is simplicity:
among the candidates that fit the data, prefer the shortest. Which candidates fit is
fixed by the record. Which of them is shortest depends on the reference machine, the
universal computer in which lengths are measured, and there are always rivals that fit
equally well, are nearly as short, and disagree about what has not been observed.

\subsection{Description length as a constraint}

\begin{definition}[Prefix description system]\label{con:def:prefix}
A \emph{prefix description system} is a partial computable function
$U:\{0,1\}^*\times\Str\rightharpoonup\Str$ whose set of programs
$D_U=\{p : U(p,x)\conv \text{ for some } x\}$ is prefix-free. A program
$p\in D_U$ determines the model $x\mapsto U(p,x)$; by the $s$-$m$-$n$ theorem an index
$e_p$ with $\ph{e_p}=U(p,\cdot)$ is computable from $p$. $U$ is \emph{universal} if
for every prefix description system $V$ there is a string $\sigma_V$ with
$U(\sigma_V q,x)\simeq V(q,x)$ for all $q,x$. Universal systems exist
\citep{chaitin1975,li-vitanyi2019}.
\end{definition}

\begin{definition}[Description length of a record]\label{con:def:KR}
For a universal $U$ and a nonempty finite functional record $\rec$, a program
$p\in D_U$ is \emph{consistent with $\rec$} if $U(p,x)=y$ for all $(x,y)\in\rec$.
Let
\[
  \KU{U}(\rec)=\min\{|p| : p\in D_U \text{ consistent with } \rec\},
\]
and let the \emph{MDL selection relative to $U$} be the set of models $e_p$ for which
$p$ is consistent with $\rec$ and $|p|=\KU{U}(\rec)$. For a string $s$ we write
$\KU{U}(s)=\KU{U}(\{(\varepsilon,s)\})$; $\K$ with no subscript denotes $\KU{U}$ for a
fixed reference $U$.
\end{definition}

$\KU{U}(\rec)$ is finite, because a lookup table for $\rec$ is a consistent program.
Minimum description length (MDL) is a constraint in the sense of
\cref{def:constraint}: let $c_U(e)=\min\{|p| : p\in D_U,\ e_p=e\}$, with
$c_U(e)=\infty$ if there is no such $p$. Since $p$ is consistent with $\rec$ exactly
when $e_p\in\Cons(\rec)$, $\Sel_{c_U}(\rec)$ is the MDL selection relative to $U$. It is
\emph{imposed} in the sense of \cref{def:constraint}: by \cref{thm:con-machine}(i)
below, the same record yields different MDL selections under different universal
machines.

\subsection{The choice of machine}

\begin{theorem}[T3: MDL is machine-relative]\label{thm:con-machine}
Let $U$ be universal, $\rec$ a nonempty finite functional record, and
$e\in\Cons(\rec)$. There is a universal prefix description system $U'$ such that
\begin{enumerate}[label=(\roman*)]
  \item for every record $\rec'\sqsupseteq\rec$ with $e\in\Cons(\rec')$, the
        MDL selection relative to $U'$ has a unique shortest consistent program,
        and it computes $\ph{e}$;
  \item $\KU{U'}(\rec')\le\KU{U}(\rec')+1$ for every nonempty functional record
        $\rec'$;
  \item $\KU{U}(\rec')\le\KU{U'}(\rec')+\KU{U}(e)+O(1)$ for every nonempty functional
        record $\rec'$, where $\KU{U}(e)$ is the length of a shortest $U$-program that
        outputs $e$ and the $O(1)$ depends only on $U$.
\end{enumerate}
\end{theorem}

\begin{proof}
Define $U'(0,x)\simeq\ph{e}(x)$ and $U'(1p,x)\simeq U(p,x)$, with $U'$ undefined on
all other programs. Then $D_{U'}\subseteq\{0\}\cup 1D_U$, which is prefix-free: $0$ is
incomparable with every string beginning with $1$, and $1p\sqsubseteq 1q$ forces
$p\sqsubseteq q$, hence $p=q$ for $p,q\in D_U$. $U'$ is universal with
$\sigma'_V=1\sigma_V$.

(i)~Since $e\in\Cons(\rec')$ and $\rec'\neq\emptyset$, $\ph{e}$ halts on some input,
so $0\in D_{U'}$ and $0$ is consistent with $\rec'$. The empty string is not in
$D_{U'}$, and $1\in D_{U'}$ only if $\varepsilon\in D_U$, which prefix-freeness and
universality exclude. So $0$ is the unique shortest consistent program.

(ii)~Each $U$-program $p$ becomes the $U'$-program $1p$.

(iii)~Let $V$ read a $U$-program $a$ (self-delimiting, since $a\in D_U$), decode
$e=U(a,\varepsilon)$, then read a $U'_e$-program $q$ and simulate $U'_e(q,\cdot)$, where
$U'_e$ is the machine $U'$ built from $e$; the simulation is uniform in $e$. Each
$U'$-program $q$ then has the $U$-program $\sigma_V a q$, where $a$ is a shortest
$U$-program for $e$, so $|a|=\KU{U}(e)$.
\end{proof}

The invariance theorem \citetext{\citealp[\S3]{kolmogorov1965}; \citealp{li-vitanyi2019}} says that two universal
machines assign description lengths within an additive constant of each other.
\Cref{thm:con-machine} says what that constant means on finite data. One extra bit in
the reference machine (ii) makes any chosen consistent candidate the unique simplest
one, and it stays so until the record refutes it (i). The price, measured in the
original machine, is at most about $\KU{U}(e)$, the length of the candidate being favored
(iii). Promoting a candidate far from the minimum, such as a lookup table for the
record, needs a change of machine with large translation constants (the converse
after \cref{thm:con-gap}). Promoting a near-tie needs only a
small one, in the following sense.

\begin{definition}[Pointwise cost of a change of machine]\label{con:def:pointwise}
Universal prefix description systems $U,U'$ differ \emph{pointwise by at most $c$} if
$|\KU{U}(\rec')-\KU{U'}(\rec')|\le c$ for every nonempty finite functional record
$\rec'$.
\end{definition}

The pointwise cost bounds how far a change of machine moves any description length.
It is not the translation constant of the invariance theorem, the length of a prefix
$\sigma$ with $U(\sigma q,x)\simeq U'(q,x)$ for all $q,x$; that constant bounds the
pointwise cost from one side, but it need not be small.

\begin{theorem}[T3, gap form: near-ties are machine-relative]\label{thm:con-gap}
Let $U$ be universal, $\rec$ a nonempty finite functional record, $\delta\ge0$ an
integer, and $q\in D_U$ a program consistent with $\rec$ with
$|q|\le\KU{U}(\rec)+\delta$. There is a universal prefix description system $U'$ such
that
\begin{enumerate}[label=(\roman*)]
  \item for every record $\rec'\sqsupseteq\rec$ with which $q$ is consistent, the
        unique shortest $U'$-program consistent with $\rec'$ computes the model
        $x\mapsto U(q,x)$;
  \item $\KU{U'}(\rec')\le\KU{U}(\rec')+1$ and $\KU{U}(\rec')\le\KU{U'}(\rec')+\delta$
        for every nonempty finite functional record $\rec'$, so $U$ and $U'$ differ
        pointwise by at most $\max(\delta,1)$.
\end{enumerate}
\end{theorem}

\begin{proof}
Let $\ell=\max(0,|q|-\delta-1)$ and $t=0^\ell$. Define $U'(0s,x)\simeq U(s,x)$ for
every $s$ and $U'(1t,x)\simeq U(q,x)$, with $U'$ undefined on all other programs. Then
$D_{U'}\subseteq 0D_U\cup\{1t\}$, which is prefix-free: $1t$ is incomparable with
every string beginning with $0$, and $0s\sqsubseteq 0s'$ forces $s=s'$ for
$s,s'\in D_U$. $U'$ is universal with $\sigma'_V=0\sigma_V$. Every $U'$-program other
than $1t$ has the form $0s$ with $s\in D_U$ and computes the same model as $s$, so
\begin{equation}\label{con:eq:gap}
  \KU{U'}(\rec')=
  \begin{cases}
    \min\bigl(\KU{U}(\rec')+1,\ \ell+1\bigr) & \text{if $q$ is consistent with $\rec'$},\\
    \KU{U}(\rec')+1 & \text{otherwise}.
  \end{cases}
\end{equation}
Note that $\KU{U}(\rec')\ge1$ for every nonempty $\rec'$: the empty program is not in
$D_U$, since otherwise $D_U=\{\varepsilon\}$ and $U$ would not be universal.

(i)~Let $\rec'\sqsupseteq\rec$ with $q$ consistent with $\rec'$. A program consistent
with $\rec'$ is consistent with $\rec$, so $\KU{U}(\rec')\ge\KU{U}(\rec)\ge|q|-\delta$.
Hence $\ell+1\le\max(1,|q|-\delta)\le\KU{U}(\rec')<\KU{U}(\rec')+1$. So $1t$, which is
consistent with $\rec'$, is strictly shorter than every consistent program $0s$, and it
computes $x\mapsto U(q,x)$.

(ii)~The first inequality is immediate from \eqref{con:eq:gap}. For the second, if
$\KU{U'}(\rec')=\KU{U}(\rec')+1$ there is nothing to prove. Otherwise
$\KU{U'}(\rec')=\ell+1$ and $q$ is consistent with $\rec'$, so
$\KU{U}(\rec')\le|q|\le\ell+1+\delta$.
\end{proof}

Two features of the theorem matter. First, the cost is pointwise only. The proof gives
no simulation of $U'$ in $U$ shorter than about $\KU{U}(q)$ bits, since $U'$ must be
told $q$; so the translation constant may be as large as the bound of
\cref{thm:con-machine}(iii), while no description length moves by more than
$\max(\delta,1)$. Second, the converse holds for translation constants. If
$U(\sigma q',x)\simeq U'(q',x)$ for all $q',x$ and $\KU{U'}\le\KU{U}+c'$, then the
shortest $U'$-program consistent with $\rec$, prefixed by $\sigma$, is a $U$-program
for the same model, so the $U'$-selection lies within $|\sigma|+c'$ bits of the
$U$-minimum. A change of machine with small translation constants reorders only
near-ties.

\begin{remark}[``Unnatural'' machines]\label{con:rem:natural}
The standard reply is that $U'$ is unnatural. But naturalness of the
reference machine is itself an imposed constraint, and no known principle privileges
one universal machine over another. The most direct attempt, a stationary distribution
for the process of universal machines simulating one another, does not exist on the
set of all universal machines; it exists only on restricted subclasses, which
reintroduces the choice \citep{muller2010}.
\end{remark}

\subsection{Near-minimal rivals}

\begin{proposition}[Near-minimal rivals]\label{thm:con-nearmin}
Let $U$ be universal, $\rec$ a nonempty finite functional record, and $p$ a shortest
program consistent with $\rec$. For every $x\notin\dom(\rec)$ and every $y\in\Str$
there is a program $q$ consistent with $\rec$ such that $U(q,x)=y$ and
\[
  |q|\le\KU{U}(\rec)+\K(x,y)+O(1),
\]
where the $O(1)$ depends only on $U$. In particular, for every such $x$ there is a
consistent $q$ that disagrees at $x$ with the shortest program $p$ (that is,
$U(q,x)\ne U(p,x)$, including the case $U(p,x)\divg$) and
$|q|\le\KU{U}(\rec)+\K(x)+O(1)$.
\end{proposition}

\begin{proof}
Let $V$ be the prefix description system that, on program $s$ and input $z$, finds the
unique prefix $a$ of $s$ in $D_U$ (by dovetailing $U$ on the prefixes of $s$), requires
$U(a,\varepsilon)=\langle x,y\rangle$, confirms that the remainder $p$ of $s$ lies in
$D_U$, and then outputs $y$ if $z=x$ and $U(p,z)$ otherwise. $D_V\subseteq\{ap : a,p\in
D_U\}$, which is prefix-free. Take $a$ a shortest $U$-program for $\langle x,y\rangle$
and set $q=\sigma_V ap$. For $(x',y')\in\rec$ we have $x'\ne x$, so
$U(q,x')=U(p,x')=y'$; and $U(q,x)=y$. For the last claim, since $|\Sig|\ge2$, choose a
single symbol $y$ different from $U(p,x)$ (any symbol if $U(p,x)\divg$); then
$\K(x,y)\le\K(x)+O(1)$.
\end{proof}

\begin{corollary}[Near-minimal rivals can be promoted]\label{con:cor:promote}
Let $U$, $\rec$ and $p$ be as in \cref{thm:con-nearmin}, and $x\notin\dom(\rec)$. There
is a universal $U'$ that differs from $U$ pointwise by at most $\K(x)+O(1)$ and whose
unique shortest program consistent with $\rec$ computes a model that disagrees with $p$
at $x$.
\end{corollary}

\begin{proof}
Take the rival $q$ of \cref{thm:con-nearmin} that disagrees with $p$ at $x$, and apply
\cref{thm:con-gap} with $\delta=|q|-\KU{U}(\rec)\le\K(x)+O(1)$.
\end{proof}

Two margins are in play here, and they are different. The first is the margin of a
\emph{change of machine}: by \cref{con:cor:promote}, a change that moves no description
length by more than $\K(x)+O(1)$ bits reverses which candidate MDL selects, in favor of
one that predicts differently at $x$. For a setting $x$ of bounded complexity this
margin is a constant. For the canonical record of length $n$ every unobserved setting
is some $x_m$ with $m\ge n$; the next one costs $\K(x_n)\le\K(n)+O(1)$, at most about $\log n$
bits, and some later ones, such as $x_{2^k}$ with $2^k\ge n$, cost only about
$\log\log n$.

The second margin is that of \emph{algorithmic statistics} for a fixed machine.
Algorithmic statistics draws the most that can be drawn from a record without a choice
among near-ties. The Kolmogorov structure function and the minimal sufficient statistic
split the record's information into a meaningful part, a simplest model for which the
record is typical, and an accidental remainder
\citep{gacs2001,vereshchagin-vitanyi2004}. Three features of that split leave the
near-ties standing. It is determined only up to additive terms logarithmic in the
length $n$ of the record; its models are sets or distributions that contain the
record, so it describes what has been observed, not how the record extends; and it is
neither computable nor usefully approximable \citep[Sec.~VII]{vereshchagin-vitanyi2004}.
The rivals of \cref{thm:con-nearmin} at the next setting cost at most
$\K(n)+O(1)$, about $\log n$, extra bits, so they lie inside the margin within which
the structure function, for any fixed machine, declares itself indifferent. This margin
does not depend on a change of machine.

\subsection{Other preferences, and serious rivals}

Simplicity is not the only preference at work in the search for a final theory.
Naturalness, beauty, symmetry, and unification are also orderings on candidate models,
and so constraints in the sense of \cref{def:constraint}: each ranks models that the
record cannot tell apart.

Nor are the surviving rivals only gerrymanders. Physics supplies rivals of comparable
compactness that agree on every prediction. When two formulations are related by a
computable, prediction-preserving translation $\tau$, their description lengths differ
by at most $\K(\tau)+O(1)$: a program for one formulation, prefixed by a
self-delimiting program for $\tau$, is a program for the other. For the explicit
translations below $\K(\tau)$ is the length of a short procedure, so if one formulation
is a shortest consistent program, the other is a near-tie of at most $\K(\tau)+O(1)$
bits, and by \cref{thm:con-gap} a change of machine that moves no description length by
more than that makes it the unique selection. Newtonian, Lagrangian and Hamiltonian
mechanics agree, for conservative systems with a regular Lagrangian, through the
Legendre transform, while their primitives differ: forces along trajectories, a
variational principle over histories, flows on phase space. Heisenberg's matrix
mechanics and Schr\"odinger's wave mechanics began from incompatible pictures and were
shown equivalent, incompletely by Schr\"odinger \citep{schrodinger1926equiv} and
rigorously by von Neumann \citep{vonneumann1932}; string-theoretic dualities are
candidates of the same kind (\cref{sec:ideal}).

\section{What even an ideal agent cannot know}\label{sec:ideal}

An agent with unlimited computing power escapes \cref{thm:ud-gold}. In a deterministic
computable world every divergent rival is refuted after finitely many observations
(\cref{ud:rem-transient}), and enumeration with an oracle for totality converges on the
true behavior (\cref{sec:discover}). Two things remain even for such an agent: rival
\emph{formulations} of one final theory, and rival \emph{chance hypotheses}, which leave
the final theory itself undetermined.

\begin{theorem}[T5: Permanent underdetermination]\label{thm:t5-permanent}
\leavevmode
\begin{enumerate}[label=(\alph*)]
  \item \emph{Formulations.} If two models are extensionally equal
        (\cref{def:model})---deterministic models $e,e'$ with $\ph{e}=\ph{e'}$, or
        stochastic models $i,j$ with $\nu_i=\nu_j$---then they assign the same
        outcome, or the same probability, to every possible record. No method of any kind, computable or not, separates
        them.
  \item \emph{Chance hypotheses.} Let $\mu$ and $\mu'$ be distinct computable
        measures on $\Seq$ that are mutually absolutely continuous, that is, agree on
        which events have probability zero: $\mu(B)=0$ if and only if $\mu'(B)=0$, for
        every Borel set $B$. Let $L$ be any method, computable or not, mapping finite outcome
        strings to models. If $L$ converges to a model of $\mu$ with $\mu$-probability
        $1$, then it converges to a model of $\mu'$ with $\mu'$-probability $0$. In
        particular, no method identifies the correct one of $\mu,\mu'$ with probability
        $1$ under both.
\end{enumerate}
\end{theorem}

\begin{proof}
(a) A record's outcomes are fixed by its settings and the model's behavior on them
(\cref{def:model}); extensionally equal models agree on both, so every method receives
the same input whichever of them is true.

(b) For a set $I$ of models (a countable set of indices), the event that $L$ converges
to a member of $I$ is
\[
  A_I=\bigcup_{N}\,\bigcup_{i\in I}\,\bigcap_{n\ge N}\{\omega\in\Seq : L(\omega_{\le n})=i\}.
\]
Each set $\{\omega : L(\omega_{\le n})=i\}$ depends only on the first $n$ symbols, so it
is a finite union of cylinders; the unions and intersections are countable. Hence $A_I$
is Borel whatever $L$ is. Let $I$ and $I'$ be the sets of models of
$\mu$ and of $\mu'$, that is, of indices $i$ with $\nu_i=\mu$ and with $\nu_i=\mu'$. They are disjoint because $\mu\neq\mu'$, so $A_I\cap A_{I'}=\emptyset$.
If $\mu(A_I)=1$ then $\mu(A_{I'})=0$, and by absolute continuity $\mu'(A_{I'})=0$.
\end{proof}

The argument uses no property of $\mu,\mu'$ beyond mutual absolute continuity. The
restriction to computable measures loses nothing: a probability measure that has a
stochastic model is lower semicomputable, and a lower semicomputable probability
measure is computable, since $\nu(r)=1-\sum_{r'\ne r,\,|r'|=|r|}\nu(r')$ is also upper
semicomputable \citep{li-vitanyi2019}.

A method that uses internal randomness is covered as well, provided each conjecture
depends on $\omega_{\le n}$ and on finitely many of the method's random bits $\rho$
(possibly a number of them that depends on the bits read). Then the success event $B$
is a Borel subset of $\Seq\times\{0,1\}^{\infty}$, and (b) applies to the products
$\mu\times\lambda$ and $\mu'\times\lambda$, where $\lambda$ is the uniform measure on the
random bits. These products are again mutually absolutely continuous: by Fubini's
theorem, $(\mu\times\lambda)(B)=\int\mu(B^{\rho})\,d\lambda(\rho)$, which vanishes if and
only if $\mu(B^{\rho})=0$ for $\lambda$-almost every $\rho$, hence if and only if
$\mu'(B^{\rho})=0$ for $\lambda$-almost every $\rho$, hence if and only if
$(\mu'\times\lambda)(B)=0$.

\paragraph{Instances of (a).}
The trivial instance, as an illustration only, is the padding lemma: every behavior has
infinitely many programs (\cref{thm:ud-infinite}(b)). The serious instances are rival \emph{formulations} that
physicists take seriously.
\begin{itemize}
  \item \textbf{Dualities.} Pairs of theories with very different ingredients are held to
        describe the same physics. In T-duality, string theory on a circle of radius
        $R$ is equivalent to string theory on a circle of radius $\alpha'/R$
        \citep{polchinski1998}. In the AdS/CFT correspondence, string theory on an
        anti-de Sitter space is conjectured to be equivalent to a gauge theory without
        gravity on its boundary \citep{maldacena1998}. Where such a duality holds, no
        record says which side is fundamental, or whether the space has one size or the
        other.
  \item \textbf{Interpretations of quantum mechanics.} The following all assign Born-rule
        probabilities to every record:
        \begin{itemize}
          \item Copenhagen-style quantum mechanics, in which the probabilities are
                fundamental;
          \item de Broglie--Bohm theory in quantum equilibrium, which is deterministic,
                with randomness confined to the initial positions
                \citep{bohm1952,dgz1992};
          \item Everettian quantum mechanics, which is deterministic unitary evolution,
                with Born weights for branches \citep{everett1957,wallace2012}.
        \end{itemize}
        As stochastic models they are therefore extensionally equal. The question ``is
        the world fundamentally random?'' is not settled by any record.
\end{itemize}
Two candidates do not belong on the list. Objective-collapse theories \citep{grw1986}
predict small departures from standard unitary quantum predictions (suppressed
macroscopic interference, spontaneous heating), so they are rivals that records can in
principle refute, not instances of (a). Superdeterministic proposals are programmatic:
their proponents describe the models as ``still in their infancy''
\citep[Sec.~5]{hossenfelder-palmer2020}.

\paragraph{Instances of (b).}
The simplest example is two coin models. Under $\mu$ every toss lands heads with
probability $1/2$. Under $\mu'$ the first toss lands heads with probability $1/3$ and
every later toss with probability $1/2$. These are distinct chance hypotheses, since
they disagree about a definite event. Because they differ in only finitely many
conditional probabilities, each of which is strictly between $0$ and $1$, they are
mutually absolutely continuous (the density of $\mu'$ with respect to $\mu$ is $2/3$ or
$4/3$). More generally, two theories of chance are mutually absolutely continuous if
their conditional probabilities differ at only finitely many histories in total, at
each of which both conditionals lie in $(0,1)$: the density is then a finite product of
ratios bounded away from $0$ and $\infty$. For independent trials with pairwise equivalent
factors $\mu_n,\mu'_n$, Kakutani's criterion gives mutual absolute continuity exactly
when $\prod_n\int\sqrt{d\mu_n\,d\mu'_n}>0$; for Bernoulli factors with probabilities
$p_n,q_n$ bounded away from $0$ and $1$, this holds if and only if
$\sum_n(p_n-q_n)^2<\infty$ \citep{kakutani1948}. A method may still guess well: if the
first toss has heads probability $0.99$ under one hypothesis and $0.01$ under the other,
guessing from that toss is right with probability $0.99$ either way. What no method
achieves is convergence on the correct hypothesis with probability $1$ under both.
Positive results exist for restricted classes. \Citet{vitanyi-chater2017} identify a
computable measure in the limit when the data are independent draws from a
computably enumerable class of measures; for dependent data their learner returns a
measure for which the sequence is typical, and there may be several, which is the gap
(b) exhibits. Without such restrictions, algorithmic identification of probabilities
fails \citep{bienvenu2018}.

\paragraph{Scope.}
The theorem takes a physical formulation to be a model in the sense of \cref{def:model},
so that whether a formulation is fundamental is a property of the model rather than of
its predictions. Part (a) then shows that the \emph{fundamental formulation} of a theory
of everything is never settled by data: two candidates with the same predictions may
disagree about what exists (strings in one space or a gauge theory on its boundary;
particles with trajectories or a branching wave function). Part (b) reaches the final
theory itself: for theories of chance, a rival that agrees on which events have
probability zero puts the final theory, in the sense of \cref{def:final}, beyond
identification with probability 1 under both.

\section{What the search is good for}\label{sec:predict}

The preceding sections are negative about \emph{finality}. Here we explain why the
search for compact, unified theories is still worthwhile. The reason is prediction. We
first state the strongest predictive guarantee available, then show that no usable
method has it, and finally say what compression offers in its place.

\paragraph{The ideal benchmark.}
Take stochastic models as in \cref{def:model}(b), with the enumeration
$\nu_1,\nu_2,\ldots$ fixed there and the outcomes of a record forming a string
$z=z_1\cdots z_n$, and write $\K(\nu)=\min\{\K(i):\nu_i=\nu\}$ for
the least prefix complexity of an index of $\nu$. The universal mixture
\[
  \Univ(z)=\sum_{i\ge1}2^{-\K(i)}\nu_i(z)
\]
weights every model by its compressed length and predicts with
$\Univ(a\mid z)=\Univ(za)/\Univ(z)$. Up to a multiplicative constant it is Solomonoff's
a priori probability \citep{solomonoff1978,li-vitanyi2019}.

\begin{theorem}[Relative optimality]\label{thm:pred-relative}
For every lower semicomputable semimeasure $\nu$ and every string $z=z_1\cdots z_n$,
\[
  \sum_{t=1}^{n}-\log_2\Univ(z_t\mid z_{<t})\;\le\;-\log_2\nu(z)+\K(\nu).
\]
\end{theorem}

\begin{proof}
If $\nu(z)=0$ the bound is trivial. Otherwise, for every index $i$ of $\nu$ the mixture
contains the term $2^{-\K(i)}\nu(z)$, so $\Univ(z)\ge2^{-\K(\nu)}\nu(z)$
\citep[Ch.~4]{li-vitanyi2019}. The conditional
log-losses telescope to $-\log_2\Univ(z)+\log_2\Univ(\varepsilon)\le-\log_2\Univ(z)$.
\end{proof}

The bound holds for every sequence, however it was produced, and it assumes nothing
about the source. Its content is comparative: on any data, the mixture loses at most
$\K(\nu)$ bits more than any given model $\nu$. The constant depends on the model and
on the reference machine.

\begin{proposition}[No computable predictor dominates]\label{pred:prop:no-computable}
Let $\rho$ be a predictor with computable conditionals satisfying
$\sum_a\rho(a\mid z)\le1$. Write $\K(\rho)$ for the least prefix complexity of a program computing $\rho$'s
conditionals. There is a computable sequence $\omega$ on which $\rho$'s
cumulative log-loss after $n$ symbols is at least $n\log_2|\Sig|-O(1)$. Meanwhile the
deterministic model of $\omega$ has log-loss $0$ and complexity at most
$\K(\rho)+O(1)$. Hence no computable predictor satisfies a bound of the form of
\cref{thm:pred-relative}.
\end{proposition}

\begin{proof}
At each step, approximate $\rho(a\mid\omega_{<t})$ for every $a\in\Sig$ to within
$2^{-t}/|\Sig|$ and let $\omega_t$ be the symbol with the least approximation, breaking
ties in a fixed order. The least conditional is at most $1/|\Sig|$, so
$\rho(\omega_t\mid\omega_{<t})\le(1+2^{-t+1})/|\Sig|$. Summing $-\log_2$ of these gives
the bound, since $\sum_t\log_2(1+2^{-t+1})<\infty$. The construction is uniform in a
program for $\rho$.
\end{proof}

By \cref{prem:effective}, then, the benchmark is not a usable method: every usable
predictor is computable, and for each there is a computable sequence on which it does as
badly as possible. What \cref{thm:pred-relative} supplies is a sense in which weighting
models by their compressed length is an optimal bias up to $K(\nu)$: relative to every lower
semicomputable model, on every sequence, it loses at most a constant. It is an ideal;
usable approximations aim at it, and each is defeated on some computable sequence. Martins draws a related moral for physics from normative
inductive reasoning: theories should be weighed, not chosen and discarded
\citep{martins2017}.

\paragraph{What compression offers.}
The benchmark nevertheless shows what compression is good for. Every model, a lookup
table included, has outputs at settings not yet observed; what distinguishes a short
model is its weight. The universal mixture gives a model of length $\K(\nu)$ weight at
least $2^{-\K(\nu)}$, so by \cref{thm:pred-relative} predicting as that model does costs
at most $\K(\nu)$ bits, while a lookup table pays about the length of the record.
Compressing what is known therefore commits a theory cheaply to predictions about what
is not yet known, and when those predictions are tested, the bias that chose the theory
is put at risk.
\begin{itemize}
  \item Maxwell's equations unified the known laws of electricity and magnetism. They
        then predicted electromagnetic waves traveling at the speed of light, which
        Hertz produced.
  \item Dirac's equation reconciled quantum mechanics with special relativity for the
        electron. It predicted a particle of the same mass and opposite charge, and the
        positron was found four years later.
  \item Compression can also fail. Le Verrier explained the anomalous precession of
        Mercury, within Newtonian mechanics, by an unseen inner planet, Vulcan. The
        planet does not exist, and the anomaly was explained by general relativity.
\end{itemize}
This is the value of unification, and it is why the search for compact theories of
everything is worth pursuing: a more compact theory of more phenomena makes cheaper
commitments beyond the record, and so it can be tested more severely.

Predictive success does not certify finality. Every success is a finite record, and by
\cref{thm:k1-certify} no finite record certifies that a model is final. Newtonian
mechanics was confirmed by predictions for two centuries, including the discovery of
Neptune, and it was not final. A theory of everything confirmed by every test anyone
performs would be in the same position.

\section{Objections and replies}\label{sec:objections}

We answer the ten objections that bear most directly on the main claim.

\paragraph{1. ``This is just Hume, or Popper.''}
Our first result,
\cref{thm:k1-certify}, is a precise version of the Humean and Popperian point recalled
in \cref{sec:intro}, and we claim no novelty
for it. The paper adds three things the Humean and Popperian traditions do not
supply. First, a usable method need not \emph{verify} a final theory to succeed: it
could simply settle on one and stay there. \Cref{thm:ud-gold} and
\cref{ud:rem-randomized} show that no usable method is guaranteed, even with any fixed
probability, to do that on every computable world compatible with what has been seen.
Any computable implementation of Popper's falsificationism is such a method, and it
inherits the limit. Second, some underdetermination is permanent (objection~3); Popper's
asymmetry between refutation and verification says nothing about rivals that can
never be refuted. Third, the usual escape from underdetermination, preferring the
simplest survivor, depends on the choice of universal machine
(\cref{thm:con-machine}). The Humean point is the first step of the argument, not
its conclusion.

\paragraph{2. ``A method need only work in our world, not in all worlds.''}
Every method
succeeds on some class of worlds and fails on others; we can grant that a method is
reliable in ours. The question is how anyone could \emph{know} that our world lies
in its class. By the prefix form of \cref{thm:ud-gold}, every finite record---the
whole of what anyone has observed, at any time---is consistent with computable worlds
in which a usable deterministic method never settles (for a randomized method, worlds
in which its chance of settling is as small as one likes). And for any method whose
success class meets every record, such as one that identifies the eventually-zero
worlds, the same record is also consistent with worlds in which it succeeds
(\cref{lem:bias}).
The judgment that ours is one of the good worlds is therefore not drawn from the
records. It is the method's bias, restated as a belief about the world. Belot argues that Bayesian
convergence theorems make a similar confidence rationally mandatory \citep{belot2013}.

\paragraph{3. ``Underdetermination is only transient.''}
Each rival in \cref{thm:ud-infinite} disagrees with the others somewhere, so each is
refuted once the record reaches that setting. Sklar calls this ``transient''
underdetermination \citep[p.~380]{sklar1975}, and the objection is that it is only a temporary
state of the evidence;
Laudan and Leplin add that empirical equivalence is relative to a time and to
auxiliary assumptions, and may later be broken \citep{laudan-leplin1991}. We grant
both points rival by rival. Two things remain. First, at every stage infinitely
many unrefuted rivals still disagree about what comes next, and by the prefix form
of \cref{thm:ud-gold} no usable method can tell when the rivals that matter are
gone: the transient phase never announces its end. Stanford argues that recurring
transient underdetermination, in the form of alternatives no one has thought of, is
the real threat to realism \citep{stanford2006}; the rivals of \cref{thm:ud-infinite} show that unrefuted alternatives always exist, though they
are not Stanford's unconceived theories. Second, some underdetermination is not transient at all. Models that
agree on every possible observation are never separated, and chance hypotheses that
agree on which events have probability zero are never told apart with certainty by
any method (\cref{thm:t5-permanent}).

\paragraph{4. ``The rivals are cheap.''}
Kukla gave algorithms that turn any theory into an empirically equivalent rival
(\cref{sec:certify}) \citep{kukla1993}. Leplin
and Laudan reply that such constructions are not genuine scientific rivals
\citep{leplin-laudan1993}; Kukla later defended them \citep{kukla1996}; and Stanford that rivals produced by a general recipe
amount to Cartesian skepticism \citep{stanford2001}. Three replies. First, the
criteria that exclude cheap rivals---independence, naturalness, explanatory
depth---are exactly the kind of preference \cref{lem:bias} identifies: real, often
useful, and not supplied by the records. Second, the rivals need not be cheap.
\Cref{thm:con-nearmin} gives rivals within $\K(x)+O(1)$ bits of the shortest model that
predict a different outcome at any unobserved setting, so the excluding criterion
cannot be length. Third, physics supplies rivals no one regards as artificial:
Newtonian, Lagrangian and Hamiltonian mechanics; matrix and wave mechanics; dual
descriptions of the same physics; and the interpretations of quantum mechanics
(\cref{sec:simplicity,sec:ideal}). Nor is the argument skeptical: it posits no
deceiver and holds for any record, however it was produced.

\paragraph{5. ``Simplicity, or beauty, tracks the truth.''}
The hope for a final theory is usually a hope that the true theory is simple and
beautiful \citep{weinberg1992}. Five points.
\begin{itemize}
  \item \emph{Description length.} The most precise notion of simplicity depends on
        the choice of universal machine (\cref{thm:con-machine}): on any finite record,
        any consistent model can be made the shortest, at a cost of about $\K(e)$
        bits, so the record leaves open the order of candidates within that margin.
        Description length is called objective because it is invariant up to an
        additive constant. We agree, both asymptotically and for comparisons separated
        by more than that constant; but near-tied candidates that disagree about the
        unobserved always exist, and a change of machine that moves no description
        length by more than about the gap makes any of them the shortest
        (\cref{thm:con-nearmin,thm:con-gap}).
  \item \emph{Ideal MDL.} \citet{vitanyi-li2000} show that ideal MDL is Bayes-optimal
        when the data are random for the true hypothesis. We accept this. It is a
        conditional guarantee whose condition, typicality, no finite record can confirm,
        since randomness deficiency is uncomputable; and it yields convergence in the
        limit, never certification.
  \item \emph{Hard to vary.} Deutsch holds that good explanations are those whose
        details cannot be varied without spoiling their fit, and that choosing them is
        rational, not a matter of taste \citep[ch.~1]{deutsch2011}. Taken informally, this is
        one more constraint in the sense of \cref{lem:bias}; formalized as a count of
        free parameters or a description length, it inherits machine-relativity. But
        Deutsch also denies that any explanation is final: ``there can be no such thing
        as an ultimate explanation'' \citep[ch.~3]{deutsch2011}. His view does not oppose the main
        claim. He opposes us only on whether the selection is rational,
        and nothing in our argument requires that it be irrational, only that the
        records do not make it.
  \item \emph{Ockham efficiency.} Kelly and Genin prove that an Ockham method, which
        prefers the simplest answer to a fixed question, converges to the truth with
        worst-case optimal retractions \citep{kelly2007}, or optimal reversals and
        cycles of opinion \citep{genin-kelly2019}. This
        justification is not machine-relative, and we accept it. It is relative to a
        chosen question and its partition into answers, and it guarantees efficient
        convergence in the limit, never certification: an Ockham method never knows
        that it has stopped retracting. It is consistent with our claim.
  \item \emph{Beauty and naturalness.} Our point is independent of Hossenfelder's
        diagnosis (\cref{sec:intro}):
        such criteria are constraints the records do not supply, and no
        record can confirm that the world satisfies them.
\end{itemize}

\paragraph{6. ``The success of physics shows convergence on the truth.''}
The no-miracles argument holds that the predictive success of mature theories would
be a miracle unless they were approximately true \citep{putnam1975}. We accept the
success and explain it differently. A compact model of what is known commits itself
about what is not, and when those commitments survive test the model has earned its
place (\cref{sec:predict}). The ideal predictor of \cref{thm:pred-relative} enters only as a
benchmark: it shows that weighting models by compressed length is a
bias that loses little against any computable model, but no scientist uses it, and
it does not explain why particular theories succeed. The
success of physics is evidence that its theories are good compressions; it cannot
certify that any of them is the last (\cref{thm:k1-certify,sec:predict}).

\paragraph{7. ``Algorithmic probability is itself the law.''}
Müller's algorithmic idealism \citep{muller2026} shares our starting point and our tools: it begins from
the observer's record rather than from a world, and it is formalized within
algorithmic information theory \citep{muller2020}. It also anticipates \cref{thm:con-machine}. In
\emph{Law without law} \citep{muller2020} Müller holds, citing his own result that the most natural attempt to
privilege a universal machine fails \citep{muller2010}, that the choice of algorithmic prior is
arbitrary, and he restricts the theory's predictions to those that are identical for
every choice. From these postulates he derives that observers will, with high
probability, come to see an external world governed by simple computable laws. The
step we decline is his postulate that algorithmic probability \emph{gives the chances}
of what an observer sees next, which makes it the fundamental law. Read as a candidate
final theory, this postulate is subject to three limits. No
record certifies it: every finite record has positive probability under infinitely
many computable measures that disagree about what comes next
(\cref{thm:ud-infinite}(c)), so by part (b) of \cref{thm:k1-certify} any certification
of the universal mixture at a record is wrong in some computable world that produces
that record.
Adopting it is a choice of bias (\cref{lem:bias}), and Müller's own restriction to
machine-invariant predictions concedes that single-step chances are not fixed by the
theory; its machine-invariant predictions are asymptotic. And the mixture is uncomputable, so no usable agent can apply the law
it would be (\cref{pred:prop:no-computable}). Algorithmic idealism is therefore not a
counterexample to our claim but an instance of it: a constraint,
explicitly chosen, that no observation could confirm as final. The two views agree on
what records can and cannot settle, and differ only on whether to affirm a law beyond
them. Müller affirms one; we remain agnostic.

\paragraph{8. ``The second premise begs the question.''}
Refuting \cref{prem:effective} requires a \emph{usable} hypercomputer: a system whose
non-computable behavior can be specified, relied on and put to work, not one whose
user cannot rule out spurious signals or device failure \citep{earman-norton1993}.
For deterministic methods C2 survives any hypercomputer
with the power of a Turing oracle, provided theories can use the same power, because
the diagonal argument reruns at every level of computing power
(\cref{thm:ud-gold-rel,rem:p2star}). This also answers Faizal and colleagues (\cref{sec:intro}): non-algorithmic
understanding would help only if it were
unavailable to the theories it was used to choose among.

\paragraph{9. ``Knowledge needs only reliability, not proof.''}
On fallibilist and externalist accounts, a belief can be knowledge because it was
reliably produced or is safe, without the record ruling out every rival. The
externalist may say that a reliable method that has settled on the final theory, in a
world where it is final, yields knowledge in their sense. We need not deny this: a
believer in that world may hold the truth. What no one can do is establish, from the
records, that they are in that world (C1, C2, \cref{lem:bias}), and that is what
``known to be final'' means here (\cref{def:known}).

\paragraph{10. ``The argument undermines itself.''}
\Cref{prem:record,prem:effective} are universal empirical claims, so, as for any
universal claim, no record certifies them either. We do not claim to know them. We
assert them as premises and state what would refute each (\cref{sec:framework,rem:p2star};
objection~8), so every conclusion is conditional on them. This is consistent: the paper claims no certification of a kind it
shows to be impossible.

\section{Conclusion}\label{sec:conclusion}

We asked whether a theory of everything could ever be known to be final, and argued that
it could not.
\begin{enumerate}[label=(C\arabic*)]
  \item No method of any kind can certify a theory as final: if a method certifies a
        model, deterministic or stochastic, at a record, a computable world that produces
        that record makes the certification wrong (\cref{thm:k1-certify}).
  \item No usable method is guaranteed, even with any fixed probability, to settle on a
        final theory. Whatever has been observed so far, there is a possible world that
        fits all of it in which a deterministic usable method never settles on the right
        theory (\cref{thm:ud-gold}), and for a randomized one and each $p>0$, there is a world in which it settles with
        probability less than $p$ (\cref{ud:rem-randomized}). This part rests on the second premise that usable methods
        are computable (\cref{prem:effective}), and for deterministic methods only on its
        weaker form (\cref{thm:ud-gold-rel}).
  \item Even an ideal agent cannot separate formulations with identical predictions,
        and cannot tell apart with certainty chance hypotheses that agree on which events
        have probability zero (\cref{thm:t5-permanent}).
\end{enumerate}
As a corollary (C$'$), simplicity narrows the field but cannot decide among the survivors
that matter. The data fix which models fit; which of them is shortest depends on the
reference machine. Rivals that fit equally well, are nearly as short and disagree about
the unobserved always survive, and a change of machine that moves no description length by more than about a rival's gap makes it the
unique shortest (\cref{thm:con-gap,con:cor:promote}). The universal mixture still predicts
within $\K(\nu)$ bits of every model $\nu$ (\cref{thm:pred-relative}). Every choice among them is a
constraint we impose (\cref{lem:bias}).

The position is agnostic in Huxley's sense. It does not deny that the world has a true
description, or that physicists might one day write one down; it denies that anyone
could know, from the records, that they had. Nor does it counsel giving up on unification. Compact theories remain
worth seeking, because they compress what is known and commit to new, testable
predictions (\cref{sec:predict}). What the argument removes is the goal of a final theory known to be final, and with it the hope that simplicity, beauty
or unification could certify one.

Put plainly: we can keep finding better theories of everything, but we can never know,
from the records, that we have found the last one.

\section*{Statement on the use of AI}
\addcontentsline{toc}{section}{Statement on the use of AI}

This paper was written with substantial assistance from AI systems: Claude Code
(Anthropic; Claude Opus~5.5) and, in developing the initial idea, Grok (xAI). They were used to
develop and review the argument and proofs, to research and check the literature, and to draft
the text. AI systems are not authors. The author directed the work and is responsible for its
content. A record of the process is in the project repository.


\bibliographystyle{plainnat}
{\raggedright
\bibliography{refs}
}

\end{document}
